\documentclass[../main.tex]{subfiles}
\usepackage{tcolorbox}


\graphicspath{{\subfix{../}}}
\begin{document}


\section{Genetic algorithm}
\subsection{Outline}
For the second assignment, a \textbf{genetic algorithm} has been designed to solve the Traveling Salesperson Problem (TSP).

A \textit{graph} is represented using an adjacency matrix, and an hamiltonian \textit{tour} is expressed as a vector of non-repeating integer numbers, each corresponding to a node. The length of the vector corresponding to a tour matches the number of nodes in the graph. It is worthy to note that any hamiltonian tour can be expressed in many ways as a vectors of node, since a cycle has no starting node.

The hamiltonian tour is considered as the individual in our problem, and the fitness function is the inverse total weight of the edges traversed in the tour. Each individual is a valid possible solution for our problem. Although we are aware that is not a compulsory characteristic, this feature avoids us the process of "fixing" an unfeasible solution.

Various selection, crossover, and mutation operators have been implemented within a common algorithmic framework, including:

\begin{itemize}
    \item Tournament selection
    \item Ranking selection
    \item Order crossover (OX)
    \item Partially mapped crossover (PMX)
    \item Reverse Sequence Mutation (PSM)
    \item Partial Shuffle Mutation (RSM)
\end{itemize}


\subsection{Initial population generation}
The initial population consists of randomly generated permutation of nodes, each having as many node as the graph of the problem.

Few individuals (5\%) are generated using one of the following heuristics.

\subsubsection{Nearest neighbor heuristics}

Begin at a random node  and build a tour by always moving to the nearest neighbor (among the nodes not yet connected to  
the tour) of the latest node to join the tour. When all nodes have been connected  to the tour, join the last node to the origin to complete the tour.\cite{mitlogisticaltransportation}


Worst-case performance: $\frac{1}{2} \lceil \log_2 n \rceil + \frac{1}{2}$

Complexity: $O(n^2)$.

\subsubsection{Farthest insertion}
The tour construction process begins with an initial loop consisting of only two nodes. At each iteration, a new node is inserted into the tour.  

The insertion cost is defined as:  
\begin{equation}
    c(i, k, j) = d(i, k) + d(k, j) - d(i, j),
\end{equation}
where $i$ and $j$ are consecutive nodes in the current tour $T$, and $d(a, b)$ represents the distance between nodes $a$ and $b$.  

The next node to be added to the tour is the one that maximizes $d(k, T)$, i.e., the node farthest from the existing tour. Once selected, the node is inserted at the position $(i, j)$ that minimizes the insertion cost $c(i, k, j)$. This process is repeated iteratively until all nodes have been incorporated into the tour $T$.

Worst-case performance: unknown \cite{mitlogisticaltransportation}

Complexity: $O(n^2)$.
\subsection{Crossover operators}

Both the following crossover have been implemented.
\subsubsection{Partially Mapped Crossover (PMX)}

Partially Mapped Crossover (PMX) is a recombination operator used in permutation-based genetic algorithms, particularly for problems such as the Traveling Salesman Problem (TSP). It ensures that offspring maintain valid permutations while inheriting information from both parents.

Let $P_0$ and $P_1$ be the first and second parent chromosomes.
The algorithm is the following \cite{crossover-wikipedia}
\begin{enumerate}
    \item Randomly select two crossover points in $P_0$.
    \item Copy the segment between these two points from $P_0$ to the child chromosome at the same position.
    \item Identify genes that have not been copied by checking the corresponding segment in $P_1$, starting from the first crossover point.
    \item For each missing gene (denoted as $m$):
    \begin{itemize}
        \item Find the gene $n$ in the offspring that was copied from $P_0$ in its place.
        \item Copy $m$ to the position held by $n$ in $P_1$ if that position is unoccupied.
        \item Otherwise, move to the next missing gene and repeat the process.
    \end{itemize}
\end{enumerate}

This method preserves positional information and maintains valid permutations.

\subsubsection{Order Crossover (OX)}

Order Crossover (OX) is a recombination operator used in permutation-based genetic algorithms. It ensures that the offspring inherit the relative order of elements from both parents.


Let $P_0$ and $P_1$ be the first and second parent chromosomes.
The algorithm is the following \cite{crossover-wikipedia}:
\begin{enumerate}
    \item Randomly select two crossover points in $P_0$.
    \item Copy the segment between these points from $P_0$ to the child chromosome at the same position.
    \item Traverse $P_1$ from the second crossover point onward (wrapping around if necessary) and copy the remaining genes in the same order as they appear in $P_1$, skipping any genes already present in the child.
\end{enumerate}

This method preserves the relative order of genes while maintaining valid permutations in the offspring.


\subsection{Genetic operators}

Both the following crossover have been implemented.
%From now on, the
\subsubsection{Reverse Sequence Mutation (RSM)}

The Reverse Sequence Mutation
\cite{DBLP:journals/corr/abs-1203-3099}
reverses a sequence limited by two position $i$ and $j$ such that $0 \leq i < j \leq n-1$, where $n$ is the chromosome length (in our case the number of holes). 

\subsubsection{Partial Shuffle Mutation (PSM)}
The Partial Shuffle Mutation \cite{DBLP:journals/corr/abs-1203-3099}
swap one or more couple of elements in the chromosome. In our implementation a PSM mutation 
swaps only a couple of element.

\subsection{Selection operators}
We have implemented two selection operators: the linear ranking selection operator and the $k$-tournament selection operator.

The linear ranking selection operator assigns a rank to each individual in the population based on their fitness values, with higher-ranked individuals having a greater probability of being selected for reproduction. The $k$-tournament selection operator involves randomly selecting a subset of individuals from the population and choosing the best among them as a parent.

\subsection{Survival selection}

We choosed the $(\mu, \gamma)$ evolution strategy. In each iteration, $\lambda$ offspring are generated, from which the $\mu < \lambda$ best individuals are selected to form the next generation. Following \cite{oslobioinspired}, we set the offspring size as $\lambda = 3\mu$, which provides a good balance between exploration and selection pressure.



The offspring selection is based on the fitness rankings, not on the actual fitness values. The resulting algorithm is therefore invariant with respect to monotonic transformations of the objective function. 

An alternative evolution strategy is ($\mu$ + $\gamma$), in which the best are selected from a union of $\mu$ parents and $\lambda$ offspring. The strategy is considered worse than ($\mu$, $\gamma$) in leaving local optima \cite{oslobioinspired}.
\subsection{Elitism}
We let the user specify an elitism rate in [0,1], which preserves the best solutions across generations to prevent losing high-quality candidates. Specifically, given the $elitism\_rate$ we retain $elitism\_counter$ (computed as $\mu$ * $elitism\_rate$) best parents and select the best ($\mu$ - $elitism\_counter$) offspring to form the next generation's population of size $\mu$.
\section{Implementation}
\subsection{Class Design}

The following classes have been defined. Each class has been designed with minimal responsibilities, adhering to the principle of single responsibility: This explain the relatively large number of classes in the project.

Description of trivial methods such as constructors and conversion operators are omitted.

\begin{itemize}
\item The \texttt{Node} class represent a two-dimensional point, to be used as a vertex of a \texttt{Graph}. Contains a method \texttt{double distance(const Node\&) const}, that returns the distance from another \texttt{Node}.

\item The \texttt{Graph} class represent a complete weighted graph. A \texttt{Graph} is constructed from an array of \texttt{Node}, representing its vertices. The weights of the graph's edges are stored in a (private) two-dimensional vector of \texttt{double} values.

\item
The \texttt{Tour} class represents a single candidate solution and is responsible for storing and manipulating a chromosome. 

It contains a vector for the sequence of node of a hamiltonian tour, and a double representing the tour's cost. Two methods are provided to reverse a sub-path (\texttt{reverse\_sub\_path(int from, int to)}) and to swap two indices (\texttt{swap\_indexes(int i, int j)}). The tour's cost is updated to reflect the current state of the tour after construction and any operations.



\item the \texttt{FitnessEvaluator} class calculates the cost of a \texttt{Tour} based on the total weight of the edges traversed in the order specified by the \texttt{Tour} object. 
At \texttt{FitnessEvaluator} is constructed  after a \texttt{Graph} object, which contains the adjeciency matrix that is necessary information to compute the cost value.
    
\item 
The \texttt{Solver} class  represents the core of the genetic algorithm process. The method \texttt{solve} define the main loop (where the evolution take place) and its termination condition. The loop makes polimorphic calls to methods of \texttt{\_mutation\_operator}, \texttt{\_crossover\_operator}, and \texttt{\_selection\_operator}, i.e. utilizes different genetic operator depending on which object the user passed to the constructor.

An object of class \texttt{Solver} needs the caller to set a \textbf{FitnessEvaluator}, and all the genetic operators (for Mutation, Crossover, and Selection) in order for the \texttt{solve} method to be used on an population (vector) of \textbf{Tour}

\end{itemize}

\subsubsection{Genetic operators}
Three \emph{abstract} classes are defined as being the base for genetic operators:

\begin{itemize}
\item
\texttt{MutationOperator}, having a public virtual method \texttt{mutate}, that is implemented by the two concrete subclasses, \texttt{PSMutation} and \texttt{RSMutation}.
\item 
\texttt{CrossoverOperator} having a public virtual method \texttt{crossover}, that is implemented by the two concrete subclasses, \texttt{OrderCrossover} and \texttt{PMXCrossover}.
\item
\texttt{SelectionOperator} having a public virtual method \texttt{select}, that is implemented by the two concrete subclasses, \texttt{TournamentSelection} and \texttt{RankingSelection}.
\end{itemize}
\paragraph*{Extensibility}
A user can define their own genetic operators by subclassing a proper class and subsequently instantiate a \texttt{Solver} that uses the newly defined operators.

\paragraph*{Similarity with strategy pattern}
The design of class \texttt{Solver} and genetic operators classes is loosely inspired by  
the well-known strategy pattern. In our implementation, the caller can specify the strategy by providing an instance of a specific subclass to \texttt{Solver}'s constructor. However, this strategy is immutable at runtime, meaning it cannot be altered once the object has been instantiated.

%The inspiration for strategy pattern comes from 
\subsubsection{Helper classes}
\begin{itemize}
\item \texttt{TSPGraphReader} provides a method \texttt{read(const std::string \&filepath)} that returns a vector of \texttt{Node} corresponding to point indicated by a valid file in TSPLIB-format.
The TSPLIB-format is the following:
\begin{verbatim}
%..some headers..
NODE_COORD_SECTION
1   x1   y1
2   x2   y2
%..some other points
EOF
\end{verbatim}
\item \texttt{CommandLineOptions} is responsible for parsing and handling command-line arguments.
\end{itemize}


\begin{tcolorbox}[title=Note - Iterators and access operators]
Most of our data structures provide \texttt{begin()}, \texttt{end()}, and \texttt{operator[]}. These methods:	
\begin{itemize}
\item
 ensure compatibility with STL algorithms that operates on iterator pairs (such as \texttt{std::find}, \texttt{std::sort}, and \texttt{std::accumulate});
 \item allow usage in constructs like:
  \vspace{-10pt}  
\begin{verbatim}
  for (auto it = container.begin(); it != container.end(); ++it) { ... }
\end{verbatim}
  \vspace{-10pt}
and:
  \vspace{-10pt}
\begin{verbatim}
  for (auto elem : container) { ... }
\end{verbatim}
  \vspace{-10pt}
\end{itemize}
Furthermore, \texttt{cbegin()} and \texttt{cend()} return constant iterators, and are used to explicitly indicate that the referenced elements remain unmodified.
\end{tcolorbox}


\subsection{Usage}
Available options can be displayed by running \texttt{assignment2 -h}. The ouptut is the following:
\begin{verbatim}
assignment2 [options]
options:
  -f, --filename [FILE]                  	  Name of the input TSP instance file
  
  -t, --timeout-ms [TIMEOUT]                Timeout expressed in milliseconds

  -m, --mutation-probability [PROBABILITY]  Probability that a mutation occurs 
                                             in the genetic algorithm

  -c, --crossover-rate [RATE]           Probability that two selected solutions 
                                         are mated to create a new offspring

  -e, --elitism-rate [RATE]             Elitism rate

  --mu [SIZE]                           Size of the population pool

  --lambda [SIZE]                       Size of the offspring pool before being pruned

  -k, --tournament-size [SIZE]          Size of the tournament selection

  -N, --max-gen-no-improvement [NUMBER] Maximum number of generations 
                                         without solution improvement

  -M, --max-gen [NUMBER]                Maximum number of generations
  
  -b, --best [BOOL]                  If true, runs the algorithm once with the best 	
                                         combination of genetic operators; if false, 
                                         runs it 8 times with different combinations
\end{verbatim}


Any run of \texttt{assignment1 -f FILENAME} generates: 
\begin{itemize}
\item a file having extension \texttt{.lp} which contains the description of the model as obtained by \texttt{CPXwriteprob};

\item the solution to the problem (as an ordered sequence of integers representing the nodes), the optimal cost (as an euclidean distance) and the eventual MIP gap; in \texttt{stdout}.
\end{itemize}

Any run of \texttt{assignment1 -f FILENAME} generates, for each of the 8 combinations
of genetic operator implemented, the approximate solution to the problem (as an ordered sequence of integers representing the nodes) and its cost (as an euclidean distance), in \texttt{stdout}.

\section{Further improvements}

To refine the solutions produced by genetic algorithm, we propose the incorporation of local optimization techniques on the permutation generated by our algorithm; for example the 2-opt algorithm or Tabu Search.

The 2-opt algorithm systematically improves the solution by iteratively removing two crossing segments and reconnecting the resulting paths, thereby reducing the overall tour length through incremental improvements: In facts 2-opt is guaranteed to never output a solution worse that the previous one.

Alternatively, Tabu Search can be used to explore the solution space more broadly. By maintaining a tabu list of recently visited solutions, Tabu Search allows for the acceptance of non-improving moves, which helps to escape local optima and encourages a more diverse search.

A different fitness function that penalizes solution similar to the current best one might be tried.

\end{document}
